\runninghead{The Leaven in Three Measures}
\properlane{hidden-leaven}{Leaven Hidden in Three Measures: Love at Work in the Whole Person}

The second parable compares the kingdom to something that the New Testament
elsewhere uses for corruption. ``Take heed and beware of the leaven of the
Pharisees and Sadducees,'' Jesus warns (Mt 16:6), and Paul tells the Corinthians to purge out the old
leaven and to keep the feast not with the leaven of malice and wickedness
(1~Cor 5:7--8). Here the woman takes leaven, \latin{abscóndit}, ``hid,'' it in
three measures of meal, and leaves it \latin{donec fermentátum est totum},
``until the whole was leavened.'' The Fathers ask what this leaven is, what
the three measures are, and how far ``the whole'' reaches. St Augustine and St
Bede answer that the leaven is love; St Jerome and St Ambrose, that it is the
teaching of Scripture and of the Church; and all four take the three measures
to be the whole human person. Bede adds an argument from the word
\latin{donec}: love hidden in the mind must go on growing until the whole mind
is changed. This interpretation takes the charity for which Paul gives thanks
in the Thessalonians to be love of that kind, already at work in a church.

\subsection{Leaven for good and leaven for evil}
\label{sec:leaven-difficulty}

St Thomas Aquinas states the difficulty and its answer in one distinction.
The same thing may be taken \latin{in malum} or \latin{in bonum}: leaven is
evil by its corruption, as when Paul says ``purge out the old leaven'' (1~Cor
5:7), and good by its heat and its power of spreading through the dough
(\work{Super Evangelium Matthaei} c.~13). Ambrose answers from the very texts
that raise the difficulty. The Lord's warning against the leaven of the
Pharisees and Paul's ``not with the leaven of malice and wickedness'' show,
he says, that leaven means teaching, \latin{ostenditur quia doctrina fermentum
est}; and \latin{aliud est fermentum zizaniae, aliud fermentum tritici}, ``the
leaven of the cockle is one thing, the leaven of the wheat another''
(\work{Expositio in Lucam}~VII.190).

\subsection{What the leaven is}
\label{sec:leaven-what}

Ambrose admits that the parable \latin{tantum ambiguitatis offundit, ut
diversas habeat sententias plurimorum}, ``spreads so much ambiguity that it has
the different opinions of many'' (VII.187), and he gives several before his
own. Many, he reports, take the leaven to be Christ, and on their reading holy
Church, figured by the woman, \latin{Dominum Iesum in interioribus nostrae
mentis abscondat, donec animi nostri secreta penetralia calor sapientiae
coelestis obducat}, ``hides the Lord Jesus in the inner parts of our mind,
until the warmth of heavenly wisdom covers the secret recesses of our soul''
(VII.187). Because Matthew has three measures, it seemed fitting to him also
to believe the Son of God hidden in the Law, covered in the Prophets and
fulfilled in the preaching of the Gospel (VII.188); and some, he says, take
the parable of the world (VII.189). Then he gives his own view: \latin{Me
tamen sequi iuvat quod ipse Dominus docuit fermentum Ecclesiae spiritalem esse
doctrinam}, ``Yet I am glad to follow what the Lord himself taught, that the
leaven of the Church is spiritual teaching'' (VII.190).

Jerome reaches a similar answer by a shorter road. The woman, he says,
\latin{mihi videtur}, ``seems to me,'' to be the apostolic preaching or the
Church, and the leaven is the knowledge of the Scriptures (\work{In
Matthaeum}~II, on 13:33, PL~26 col.~91). Augustine gives the other answer in
two lines: \latin{mulierem, sapientiam dicit: fermentum, dilectionem; quod
fervefacit et excitat}, ``by the woman he means wisdom; by the leaven, love,
which warms and rouses'' (\work{Quaestiones Evangeliorum}~I.12). Bede, on
Luke's form of the parable, takes up Augustine's words, adding the object:
love \latin{quae fervefacit, et excitat mentem}, ``which warms and rouses the
mind'' (\work{In Lucam}~IV, on 13:21).

Others answer differently again. St Hilary of Poitiers takes the leaven to be
the Lord, hidden by the synagogue \latin{per judicium mortis}, ``through the
sentence of death'' (\work{In Matthaeum}~13.5). St John Chrysostom takes it to
be the disciples, mixed into the multitude of the world, and St Peter
Chrysologus, in one sermon, both the leaven of resurrection that replaces the
leaven of death in the human race and \latin{caelestis doctrinae
sacramentum}, ``the mystery of heavenly teaching'' (\work{Sermo}~99).

\subsection{Three measures: the whole person}
\label{sec:leaven-measures}

Augustine, Bede, Jerome and Ambrose agree more closely on the three measures
than on the leaven. Augustine names them from the great commandment:
\latin{Ex toto corde, et ex tota anima, et ex tota mente}, ``with the whole
heart, and with the whole soul, and with the whole mind'' (Mt 22:37); he adds
two other possibilities, the hundredfold, sixtyfold and thirtyfold of the
sower, and Noe, Daniel and Job (\work{Quaestiones Evangeliorum}~I.12). Bede
makes the commandment the reason the leaven is hidden in three measures at
all: \latin{Abscondit ergo mulier fermentum dilectionis in farinae sata tria,
quia praecipit Ecclesia ut diligamus Dominum ex toto corde, tota anima, tota
virtute}, ``The woman hides the leaven of love in three measures of meal,
because the Church commands us to love the Lord with the whole heart, the
whole soul and the whole strength'' (\work{In Lucam}~IV, on 13:21).

Jerome takes the measures as spirit, soul and body, and the leaven's work as
their agreement: the Scriptures are hidden in them \latin{ut spiritus, anima,
et corpus in unum redacta, non discrepent inter se}, ``so that spirit, soul
and body, brought into one, may not be at odds among themselves.'' He gives
another exposition after Plato's three parts of the soul, and takes it up for
himself and his readers: \latin{ut in ratione possideamus prudentiam; in ira,
odium contra vitia; in desiderio, cupiditatem virtutum}, ``that in reason we
may possess prudence; in anger, hatred of vices; in desire, longing for the
virtues'' (\work{In Matthaeum}~II, PL~26 col.~91). Ambrose, agreeing, he says,
with good authors (\latin{bonis auctoribus acquiescimus}), takes the measures
as flesh, soul and spirit, the reborn man whom the Church sanctifies with
spiritual leaven: \latin{tres enim mensurae sunt, ut dixi, carnis, animae, et
spiritus} (VII.190--191).

Bede's words about the meal give the other side of the image. The woman is
the Church, and he identifies her with the woman in Luke whom the Lord has
just healed, a daughter of Abraham bent over for eighteen years (Lk
13:10--17). We are her meal, \latin{cujus farina nos sumus}, a phrase Ambrose
had used (VII.187); and Bede adds \latin{quotquot timoris et spei exercitio,
quasi mola superiore et inferiore conterimur, ut juxta Apostolum unus panis
unum corpus multi simus in Christo}, ``as many of us as are ground by the
exercise of fear and hope, as between an upper and a lower millstone, so that,
as the Apostle says, we being many may be one bread, one body in Christ''
(1~Cor 10:17).

The Fathers disagree about one further reading of the measures. Jerome
reports a third interpretation as that of others, which refers the three
measures to faith in the Father, the Son and the Holy Ghost, and judges it:
\latin{Pius quidem sensus, sed numquam parabolae et dubia aenigmatum
intelligentia potest ad auctoritatem dogmatum proficere}, ``a pious sense
indeed, but the understanding of a parable and the doubtful meaning of
riddles can never advance to the authority of doctrines.'' Hilary, recalling
that many referred the measures to the Trinity, writes \latin{nescio an hoc ita
opinari ratio permittat}, ``I do not know whether reason permits us to think
so,'' since in Father, Son and Holy Ghost all are one without need of any
leaven mixed in from outside (\work{In Matthaeum}~13.6). Chrysologus preaches
the Trinitarian reading, from Sara's three measures (\work{Sermo}~99).

\subsection{Until the whole is leavened}
\label{sec:leaven-donec}

Chrysostom, though he reads the leaven as the disciples, explains the verb
the parable uses: ``He said not, `put,' simply, but `hid','' so that the
leaven is ``not simply close, but so as to be actually mixed'' with the meal;
and ``as the leaven, though it be buried, yet is not destroyed, but by little
and little transmutes all into its own condition'' (\work{Hom.\ in Mt.}~46.2).

Bede draws an argument from \latin{donec}. \latin{Et pulchre dicit eatenus
absconditum in farina fermentum, donec fermentaretur totum, quia charitas in
nostra mente recondita eousque crescere debet, donec totam mentem in sui
perfectionem commutet, ut nihil videlicet praeter ipsam conditoris sui
dilectionem diligere, agere, recolere anima possit}, ``And he says well that
the leaven was hidden in the meal until the whole was leavened, because love
hidden in our mind must grow until it changes the whole mind into its own
perfection, so that the soul can love, do and recall nothing but the love of
its maker.'' The parable's ``until'' becomes in Bede a requirement, and he
sets its end beyond this life: \latin{Quod hic quidem inchoari, sed ibi habet perfici, ubi
cum sit Deus omnia in omnibus, omnes uno eodemque suae dilectionis igne
calefecit}, ``This is begun here, but it is to be perfected there, where, God
being all in all, he warms all with one and the same fire of his love'' (1~Cor
15:28; \work{In Lucam}~IV, on 13:21).

Ambrose describes the leavening as slow. The Church's leavening is \latin{non
tumultuarium, nec fortuitum, sed diuturna meditatione collatum; ut haec tria
unum sint}, ``not hasty, nor haphazard, but brought about by long meditation,
so that these three may be one'' (VII.191). The woman goes on mixing the
power of spiritual teaching \latin{donec fermentetur totus ille interior homo
cordis absconditus, et in gratiam panis coelestis assurgat}, ``until the whole
inner man of the heart, hidden, is leavened, and rises into the grace of
heavenly bread''; and teaching is rightly called Christ's leaven, \latin{quia
panis est Christus}, ``because Christ is bread'' (VII.191). There is one
leaven, he says, when the flesh does not lust against the spirit nor the
spirit against the flesh (VII.192).

Ambrose's teaching is infused into the whole man, \latin{toti homini \dots\
permixtus} (VII.190), until the three agree; Jerome's is hidden so that the
three may not be at odds. In Bede's sentence the Church hides love in the
three measures by commanding it: a teaching, the great commandment, is the
means, and love is what is hidden. On that order the two answers, teaching and
love, can be held together: the teaching is received first, and the love it
commands is what works through the whole person.

\subsection{Work, labour and charity}
\label{sec:leaven-epistle}

Paul remembers before God \latin{óperis fidei vestræ, et labóris, et
caritátis}, ``the work of your faith and labour and charity'' (1~Thess 1:3).
The Greek that Chrysostom expounds reads ``the labour of love,'' and he asks
the question the phrase raises: ``what labor is it to love? Merely to love is
no labor at all. But to love genuinely is great labor. For tell me, when a
thousand things are stirred up that would draw us from love, and we hold out
against them all, is it not labor?'' His example is Jason at Thessalonica,
dragged before the rulers for having received Paul, who gave security and
sent him away (Acts 17:5--9; \work{Hom.\ in 1~Thess.}~1). Theodoret, on the
same words, says that love has its sweetest enjoyment and also a surpassing
labour, because one must bear a brother's faults, his envy, anger, pride and
ingratitude (\work{Interpretatio in 1~Thess.}, on 1:3). Aquinas, reading the
Missal's Latin, attaches the labour to faith: \latin{memores fidei vestrae
operantis et laborantis. Item, charitatis}, ``mindful of your faith, working
and labouring; likewise, of your charity''; faith is not enough \latin{nisi
habeat operationem et laborem}, ``unless it has work and labour''
(\work{Super I ad Thessalonicenses} c.~1 l.~1).

Paul gives thanks for what he remembers of them: the work of faith, the labour
and the charity. If the leaven is love, as Augustine and Bede say, the charity
he names is that leaven, and the work and the labour are where it shows.

\subsection{Thought, word and deed}
\label{sec:leaven-orations}

The Collect asks that those who pray, \latin{semper rationabília
meditántes}, ``always meditating on reasonable things,'' may carry out what
pleases God \latin{et dictis \dots\ et factis}, ``in word and in deed.''
Ambrose said that the Church's leavening of flesh, soul and spirit is slow
work, brought about by long meditation (VII.191). The Collect describes the
people themselves as meditating, and always; their constant thought
may be taken as their share in that slow work, and the \latin{et \dots\ et}
as asking that it reach their speech and their action as well.

The Secret asks that the offering \latin{mundet \dots\ et rénovet},
``cleanse and renew,'' before it asks that it ``govern and protect.'' Paul's
command about the old leaven also puts cleansing before newness: purge out the
old leaven, that you may be a new lump (1~Cor 5:7), the verse Aquinas cites
for leaven taken \latin{in malum}. Read beside Paul's command, the cleansing
the Secret asks for is the purging of the old leaven.

The Postcommunion asks that those who have been fed may \latin{semper},
always, desire the things by which they truly live. Bede says that love is
begun here and perfected only where God is all in all, so in no communicant
has it finished its work; the continued desire the Postcommunion asks for is,
on that account, the same work going on.

\subsection{The chants: prayer that changes the one who prays}
\label{sec:leaven-chants}

The Introit's antiphon puts God's thought first, \latin{Ego cógito
cogitatiónes pacis}, ``I think thoughts of peace,'' and only then the people's
calling and his hearing. The Alleluia and the Offertory then sing a calling,
the cry from the depths, and Chrysostom says of that cry that prayer from the
deepest heart, even before it receives what it asks, calms the passions,
softens anger and drives out envy (\work{Expositio in Ps.}~129.1). The
parable's leaven also works hidden, from within the meal, and this
interpretation hears the cry from the depths as prayer that works in that way
on the one who prays.

The Gradual thanks God for deliverance \latin{ex affligéntibus nos}, ``from
them that afflict us.'' Aquinas, on the psalm, says that the saints are
afflicted in body by their enemies, and in spirit \latin{quando peccant illi
qui habitant apud eos}, ``when those who live among them sin'' (\work{Super
Psalmos}~43, n.~3). Among ``them that afflict us,'' then, those who sing the
Gradual may count the neighbours whose sins afflict them in spirit. Bearing
them comes close to the labour that Theodoret assigns to love, the bearing of
a brother's envy, anger, pride and ingratitude.

The Communion promises that whatever is asked in prayer, believed, will be
done. Jerome and Ambrose, commenting on the leaven and not on Mark, both go
on from the three measures made one to prayer that obtains. Jerome continues the
sentence already quoted: spirit, soul and body brought into one, \latin{sed
cum duobus et tribus convenerint, impetrent a Patre quodcumque postulaverint},
``but when two and three have agreed, may obtain from the Father whatever they
ask'' (Mt 18:19; PL~26 col.~91). Ambrose puts the full obtaining at the
resurrection, when there will remain \latin{incorrupta corporis et animae et
spiritus \dots\ communio: et tunc omnia quae petimus, impetremus}, ``an
incorrupt communion of body and soul and spirit, and then we shall obtain all
that we ask''; many, he says, take Mt 18:19 to signify this (VII.192). If
prayer obtains when the whole person agrees, as Jerome and Ambrose say, the
Communion's ``whatsoever you ask'' is spoken to such a person. Bede, on Mark,
adds a condition for those not yet perfect:
they will surely obtain the forgiveness they ask, \latin{si tamen in se
peccantibus aliis primo ipsi dimittunt}, ``provided that they first forgive
others who sin against them'' (\work{In Marcum}~III, on 11:25).

\subsection{The disagreement about the leaven}
\label{sec:leaven-alternative}

The strongest alternative within the tradition is Chrysostom's. For him the
leaven is the disciples, hidden among the nations, and the three measures mean
only ``many''; the parable is about the mission of the Church, as the mustard
seed is: ``even so shall ye convert the whole world'' (\work{Hom.\ in
Mt.}~46.2). Aquinas reports this reading, with the apostles hidden \latin{in
tribulationibus}, ``in tribulations.'' Chrysostom's reading makes the
leaven a second parable of the mission, beside the mustard seed, and it
belongs with the first interpretation. Hilary's reading differs in a second
way: for him the one who hides the leaven is the synagogue, condemning the
Lord to death, and the leaven so hidden in law, prophets and gospels makes
them one (\work{In Matthaeum}~13.5). For Ambrose, Bede and Chrysologus the one
who hides it is the Church, and for Augustine wisdom.

\begin{fourSenses}
\item[Literal.] A woman works a little leaven into a large batch of meal,
three \latin{sata}, each holding a \latin{modius} and a half in Jerome's
note, and leaves it until the whole has risen (Mt 13:33). Paul thanks God for
a church whose faith works, which labours and loves, whose hope endures, and
which turned from idols to serve the living God (1~Thess 1:3, 9).
\item[Allegorical.] The woman is wisdom (Augustine) or the Church (Jerome,
Ambrose, Bede); for Bede, the daughter of Abraham whom the Lord raised up, and
we are her meal, ground between fear and hope into one bread, one body
(1~Cor 10:17). In Ambrose's own reading the Church mixes Christ's teaching
into the reborn man until the whole inner man rises into the grace of
heavenly bread, because Christ is bread (VII.191).
\item[Moral.] Love hidden in the mind must grow until it has changed the whole
mind, so that the soul loves, does and recalls nothing but its maker's love
(Bede); reason, anger and desire become prudence, hatred of vice and longing
for virtue (Jerome); flesh, soul and spirit are brought to agree by the
Church's slow work (Ambrose, VII.191--192). Love is labour (Chrysostom), and
its labour is to bear a brother's faults (Theodoret); faith that does not work
and labour is not enough (Aquinas).
\item[Anagogical.] What is begun here is perfected where God is all in all
and warms all with the one fire of his love (Bede; 1~Cor 15:28); if the three
measures remain one leaven in this life, \latin{erit in futurum diligentibus
Christum incorrupta communio}, ``there will be in the future, for those who
love Christ, an incorrupt communion,'' and we shall be one, like him
(Ambrose, VII.194; 1~Jn 3:2); the whole is leavened \latin{adveniente
Christo}, ``when Christ comes'' (Chrysologus, \work{Sermo}~99).
\end{fourSenses}
